\quad O comportamento do \textit{Quick Sort} depende diretamente do
balanceamento das partições geradas a cada particionamento, que por sua
vez depende da escolha do pivô. Por esse motivo, diferentemente do
\textit{Merge Sort}, os casos crescente, decrescente e aleatório podem
apresentar custos diferentes. Neste trabalho, a análise considera as duas
versões implementadas: a versão 1, com pivô no primeiro elemento, e a
versão 2, com pivô dado pela média aritmética da partição.

\quad Em ambas as versões, o particionamento percorre os elementos da
partição uma única vez com os índices \texttt{i} e \texttt{j}, realizando
trabalho proporcional ao tamanho da partição. Na versão 2, o cálculo da
média também percorre a partição uma vez, o que mantém o custo de cada
particionamento proporcional a $n$.

\subsubsection{Melhor caso -- $\Theta(n\log n)$}

O melhor caso ocorre quando, a cada particionamento, o pivô divide a
partição em duas partes de tamanhos aproximadamente iguais. Nessa
situação, cada chamada gera duas subpartições com cerca de $n/2$
elementos, e o custo total pode ser representado por:

\begin{equation}
T(n) = 2T\left(\frac{n}{2}\right) + cn
\end{equation}

em que $cn$ representa o custo do particionamento, proporcional ao
tamanho da partição.

Como existem aproximadamente $\log_2 n$ níveis de divisão, cada um com
custo total proporcional a $n$, temos:

\begin{equation}
T(n) = \Theta(n\log n)
\end{equation}

Na versão 1, isso ocorre quando o primeiro elemento de cada partição
é, aproximadamente, a mediana dos seus elementos. Na versão 2, ocorre
quando a média aritmética está próxima da mediana, como em vetores cujos
valores são distribuídos de forma uniforme, mesmo que estejam ordenados.
Nesse caso, o cálculo da média apenas soma outro custo linear ao
particionamento, sem alterar a ordem de crescimento.

\subsubsection{Pior caso -- $\Theta(n^2)$}

O pior caso ocorre quando o pivô divide a partição de forma
extremamente desbalanceada, gerando uma partição com $n-1$ elementos e
outra vazia (ou com um único elemento). Nesse caso, a recorrência é:

\begin{equation}
T(n) = T(n-1) + cn
\end{equation}

Expandindo a recorrência, obtém-se:

\begin{equation}
T(n) = c\sum_{k=1}^{n} k = c\,\frac{n(n+1)}{2} = \Theta(n^2)
\end{equation}

Na versão 1, esse comportamento ocorre quando o primeiro elemento da
partição é sempre o menor ou o maior valor dela, o que acontece tanto em
vetores já ordenados de forma crescente quanto em vetores ordenados de
forma decrescente. Em vetores crescentes, o pivô é o menor elemento e
a partição da esquerda fica vazia; em vetores decrescentes, o pivô é o
maior elemento e a partição da direita fica com um único elemento.

Na versão 2, essas entradas deixam de ser o pior caso, pois a média
aritmética de valores distribuídos uniformemente tende a ficar próxima
do centro da partição. Entretanto, o pior caso ainda é possível quando os
valores são muito discrepantes, como em um vetor cujos elementos crescem
de forma exponencial. Nesse caso, a média é maior que quase todos os
elementos, e cada particionamento separa apenas o maior valor, gerando a
mesma recorrência acima. Assim, o pior caso da versão 2 depende da
distribuição dos valores, e não apenas da ordem em que eles aparecem.

\subsubsection{Caso médio -- $\Theta(n\log n)$}

No caso médio, considera-se uma distribuição aleatória dos elementos
do vetor. Nessa situação, a posição final do pivô pode ser qualquer uma
das $n$ posições com a mesma probabilidade, e a recorrência torna-se:

\begin{equation}
T(n) = \frac{1}{n}\sum_{k=0}^{n-1}\left[T(k) + T(n-k-1)\right] + cn
\end{equation}

A solução dessa recorrência é aproximadamente $2n\ln n$, o que resulta
em:

\begin{equation}
T(n) = \Theta(n\log n)
\end{equation}

Mesmo quando as partições não são perfeitamente equilibradas, divisões
com proporção constante (por exemplo, $1/10$ e $9/10$) ainda produzem
profundidade logarítmica. Por esse motivo, o caso médio se aproxima do
melhor caso. Na versão 2, como a média tende a ficar próxima da mediana
em entradas aleatórias com valores uniformes, o comportamento esperado é
equivalente, porém com o custo adicional da soma dos elementos em cada
particionamento.

\subsubsection{Resumo das complexidades}

\begin{table}[H]
    \centering
    \caption{Complexidades do algoritmo Quick Sort}
    \label{tab:complexidade-quick}
    \begin{tabular}{lcc}
        \hline
        \textbf{Caso} & \textbf{Versão 1 (primeiro elemento)} & \textbf{Versão 2 (média)} \\
        \hline
        Melhor caso & $\Theta(n\log n)$ & $\Theta(n\log n)$ \\
        Caso médio & $\Theta(n\log n)$ & $\Theta(n\log n)$ \\
        Pior caso & $\Theta(n^2)$ & $\Theta(n^2)$ \\
        \hline
    \end{tabular}
\end{table}

\quad Conclui-se que o \textit{Quick Sort} apresenta complexidade
$\Theta(n\log n)$ nos casos melhor e médio, e $\Theta(n^2)$ no pior caso,
em ambas as versões.
